\subsubsection{The Lax-Milgram Lemma}\label{subsec:lax-milgram-lemma}

The first tool that we will develop for deriving the existence theory for elliptic equations is commonly known as the Lax-Milgram lemma; though because of its importance we designate it as a theorem. The result is simply a generalization of the Riesz Representation Theorem (Theorem~\ref{thm:hilbert-riesz-representation}) to bilinear forms that need not be symmetric.

\begin{Theorem}{Lax-Milgram}
\label{thm:lax-milgram}
Let $H$ be a Hilbert space and let
\begin{equation}\label{eq:lax-milgram-bilinear-mapping-domain}
B:H\times H\to\mathbb{R}
\end{equation}
be a bilinear mapping. Suppose there exist positive constants $c_1$ and $c_2$ such that
\begin{equation}\label{eq:lax-milgram-bilinear-boundedness}
|B[x,y]|\leq c_1\|x\|_H\|y\|_H\quad\text{for all }x\text{ and }y\text{ in }H
\end{equation}
and
\begin{equation}\label{eq:lax-milgram-bilinear-coercivity}
B[x,x]\geq c_2\|x\|_H^2\quad\text{for all }x\in H.
\end{equation}
Then for every $f\in H^*$ there exists a unique $y\in H$ such that
\begin{equation}\label{eq:lax-milgram-variational-solution}
B[x,y]=f(x)\quad\text{for all }x\in H.
\end{equation}
Furthermore, there exists a constant $C$, independent of $f$, such that
\begin{equation}\label{eq:lax-milgram-solution-dual-norm-bound}
\|y\|_H\leq C\|f\|_{H^*}.
\end{equation}

\end{Theorem}

\begin{Remark}
\label{rem:coercive-form-energy-estimate}
A mapping $B$ satisfying \eqref{eq:lax-milgram-bilinear-coercivity} for some $c_2>0$ is called \textbf{coercive}. The inequality \eqref{eq:lax-milgram-bilinear-coercivity} can be thought of as an energy estimate. (The inequality says that the energy (the norm squared) can only blow up as fast as the bilinear form.)
\end{Remark}

\begin{Remark}
\label{rem:lax-milgram-nonsymmetric-significance}
Note that by \eqref{eq:lax-milgram-bilinear-boundedness} and \eqref{eq:lax-milgram-bilinear-coercivity}
\begin{equation}\label{eq:coercive-bilinear-induced-equivalent-norm}
\|x\|_B:=\sqrt{B[x,x]}
\end{equation}
is equivalent to the original norm on $H$. Furthermore, if $B$ is symmetric, i.e.,
\begin{equation}\label{eq:symmetric-bilinear-form-condition}
B[x,y]=B[y,x]\quad\text{for all }x\text{ and }y\text{ in }H,
\end{equation}
then $B[x,y]$ defines a new inner product on $H$. Thus, in this case, the Riesz Representation Theorem (Theorem~\ref{thm:hilbert-riesz-representation}) directly implies that for every $f\in H^*$ there exists a unique $y\in H$ such that \eqref{eq:lax-milgram-variational-solution} holds. Therefore, the significance of the Lax-Milgram lemma is that it does not require $B$ to be symmetric.
\end{Remark}

We now prove the Lax-Milgram lemma.

\begin{Proof}
For every fixed $y\in H$, the mapping
\begin{equation}\label{eq:lax-milgram-fixed-second-argument-functional}
H\ni x\mapsto B[x,y]\in\mathbb{R}
\end{equation}
is bounded and linear, i.e., an element of $H^*$. Thus, by the Riesz Representation Theorem (Theorem~\ref{thm:hilbert-riesz-representation}) there exists a unique $z\in H$ such that
\begin{equation}\label{eq:lax-milgram-riesz-representing-vector}
B[x,y]=(x,z)\quad\text{for all }x\in H.
\end{equation}
Since a unique $z\in H$ can be derived for each fixed $y\in H$ we can define a mapping $A:H\to H$ by
\begin{equation}\label{eq:lax-milgram-associated-operator-definition}
z=:A(y).
\end{equation}
The question of existence of a solution of \eqref{eq:lax-milgram-variational-solution} is now translated to the question of the invertibility of $A$. That is, for any $f\in H^*$ let $z\in H$ be the unique element such that
\begin{equation}\label{eq:lax-milgram-data-riesz-representative}
(x,z)=f(x)\quad\text{for all }x\in H.
\end{equation}
Then if for every $z\in H$ we can find a unique solution $y\in H$ of $A(y)=z$, then $y$ is the unique solution of
\begin{equation}\label{eq:lax-milgram-operator-equation-equivalence}
B[x,y]=(x,A(y))=(x,z)=f(x)\quad\text{for all }x\in H.
\end{equation}
We now note some basic properties of $A$.\begin{enumerate}\item $A$ is linear. To see this note that\[\begin{aligned}
(x,A(\alpha_1y_1+\alpha_2y_2))&:=B[x,\alpha_1y_1+\alpha_2y_2]\\
&=\alpha_1B[x,y_1]+\alpha_2B[x,y_2]\\
&=\alpha_1(x,A(y_1))+\alpha_2(x,A(y_2))\\
&=(x,\alpha_1A(y_1)+\alpha_2A(y_2)).
\end{aligned}\]Since this holds for arbitrary $x$, $\alpha_i$ and $y_i$ we have shown linearity.\item $A$ is bounded. Using \eqref{eq:lax-milgram-bilinear-boundedness} we get
\begin{equation}\label{eq:lax-milgram-associated-operator-norm-calculation}
\|A(y)\|^2=(A(y),A(y))=B[A(y),y]\leq c_1\|A(y)\|\|y\|.
\end{equation}
Canceling, we get
\begin{equation}\label{eq:lax-milgram-associated-operator-upper-bound}
\|A(y)\|\leq c_1\|y\|.
\end{equation}
\item The range of $A$ is dense in $H$. To see this we use \eqref{eq:lax-milgram-bilinear-coercivity} to note that if $y\in\mathcal{R}(A)^\perp$, then
\begin{equation}\label{eq:lax-milgram-associated-operator-range-dense}
c_2\|y\|^2\leq B[y,y]=(y,A(y))=0.
\end{equation}
\item $A$ is bounded below. Using \eqref{eq:lax-milgram-bilinear-coercivity} again (now for arbitrary $y\in H$) we get
\begin{equation}\label{eq:lax-milgram-associated-operator-lower-calculation}
c_2\|y\|^2\leq B[y,y]=(y,A(y))\leq\|y\|\|A(y)\|
\end{equation}
or
\begin{equation}\label{eq:lax-milgram-associated-operator-bounded-below}
\|A(y)\|\geq c_2\|y\|.
\end{equation}
\item Combining this with Problem~\ref{ex:closed-bounded-below-operator-closed-range} implies that the range of $A$, $\mathcal{R}(A)$, is closed. Since $\mathcal{R}(A)$ is dense, $A$ is surjective.\end{enumerate}It follows that $A$ is invertible, which gives us the existence of a unique solution $y$. Finally, the estimate \eqref{eq:lax-milgram-solution-dual-norm-bound} follows from the Riesz Representation Theorem (Theorem~\ref{thm:hilbert-riesz-representation}) and the fact that $A$ is bounded below.
\end{Proof}
